\documentclass{article}
\usepackage[margin=1.5cm]{geometry}
\usepackage{amsmath}
\usepackage{enumitem}

\setlist{nosep,leftmargin=*}

\begin{document}
\twocolumn

\title{Chapter 11 Warm-Up: A 2D Particle Class}
\author{Prof. Jordan C. Hanson}
\date{}

\maketitle

\section{A Moving Particle}

A particle in two dimensions has position \((x,y)\) and velocity
\((v_x,v_y)\).  For constant velocity over a time step \(\Delta t\),
\[
x_{\rm new}=x+v_x\Delta t,\qquad
y_{\rm new}=y+v_y\Delta t.
\]

For
\[
(x,y)=(1,2),\qquad (v_x,v_y)=(3,-1),
\]
find the position after \(\Delta t=0.5\):
\[
x_{\rm new}=\underline{\hspace{1.1cm}},\qquad
y_{\rm new}=\underline{\hspace{1.1cm}}.
\]

\section{Build the Class}

Create a class that stores the particle's state:

\begin{verbatim}
class Particle2D:
    def __init__(self,x,y,vx,vy):
        self.x = x
        self.y = y
        self.vx = vx
        self.vy = vy
\end{verbatim}

Create one particle:

\begin{verbatim}
p1 = Particle2D(1, 2, 3, -1)
\end{verbatim}

Identify the four instance attributes of \texttt{p1}.

\section{Change the State}

Add this method inside the class:

\begin{verbatim}
def update(self, dt):
    self.x += self.vx * dt
    self.y += self.vy * dt
\end{verbatim}

Then run

\begin{verbatim}
p1.update(0.5)
print(p1.x, p1.y)
\end{verbatim}

Does the output agree with your hand calculation?  Why do
\texttt{vx} and \texttt{vy} remain unchanged?

\section{Advance Through Time}

Reset the particle:

\begin{verbatim}
p1 = Particle2D(1, 2, 3, -1)

dt = 0.1
for i in range(10):
    p1.update(dt)
\end{verbatim}

How much total time has passed?
\[
t=\underline{\hspace{1.2cm}}
\]

Predict the final position:
\[
x=\underline{\hspace{1.2cm}},\qquad
y=\underline{\hspace{1.2cm}}.
\]

Print \texttt{p1.x} and \texttt{p1.y} to check.

\section{Independent Objects}

Create another particle:

\begin{verbatim}
p2 = Particle2D(-2, 0, 1, 2)
\end{verbatim}

Run

\begin{verbatim}
p1.update(0.5)
p2.update(0.5)
\end{verbatim}

Why does updating \texttt{p1} not change \texttt{p2}?

\medskip
Each object has its own \textbf{state}: its own values of
\texttt{x}, \texttt{y}, \texttt{vx}, and \texttt{vy}.

\section{Challenge: A New Method}

Add a method that returns
\[
v=\sqrt{v_x^2+v_y^2}.
\]

Begin the file with

\begin{verbatim}
import math
\end{verbatim}

and complete

\begin{verbatim}
def speed(self):
    return ...
\end{verbatim}

Test it using

\begin{verbatim}
print(p1.speed())
\end{verbatim}

Why should \texttt{speed()} return a value while
\texttt{update()} changes the object's state?

\end{document}
